\section{Implementation}
\label{sec:implementation}

This section presents the implementation of CloudAware-MOCCA evaluated in
Section~\ref{sec:evaluation}. The prototype comprises a Kotlin-based Android
middleware and two containerized Python services representing the edge and
cloud execution tiers. The MAPE-K control model is realized through explicit
monitoring, analysis, planning, execution, and knowledge-management components.
The deterministic and learned decision mechanisms operate as interchangeable
planning backends and share the same context acquisition, communication,
execution, fallback, and instrumentation infrastructure.


\subsection{System Architecture and Deployment Topology}
\label{subsec:implementation_architecture}

The deployment comprises three execution tiers: the mobile device, an edge
service constrained to 2~vCPU and 2~GB of memory, and a cloud service allocated
4~vCPU and 4~GB of memory. The two remote services execute as Docker containers
on a host separate from the mobile device and communicate through a common
container network. This topology enables transparent forwarding from the edge
tier to the cloud tier without requiring additional coordination by the mobile
application.

On the Android device, the middleware initializes the context-acquisition
components, MAPE-K controller, deterministic and learned decision policies,
connection manager, remote-execution client, and metrics recorder. These
components are maintained by a foreground Android service so that context
monitoring and task-level adaptation remain operational independently of the
user-interface lifecycle.

The edge and cloud tiers are implemented as structurally similar FastAPI
services exposing health, offloading, and status interfaces. Both services
dispatch requests through a common task registry, thereby avoiding independent
implementations of the same workload. Their principal difference concerns
resource governance. The edge service monitors container-level processor and
memory utilization and can forward requests to the cloud service when
predefined utilization thresholds are exceeded. The cloud service is
provisioned with greater computational resources and therefore represents the
less constrained remote tier.

\begin{figure}[t]
\centering
\includegraphics[
    width=\textwidth,
    height=0.35\textheight,
    keepaspectratio
]{MOCCA_Architecture.png}
\caption{CloudAware-MOCCA architecture and principal execution paths.}
\label{fig:architecture}
\end{figure}

As illustrated in Fig.~\ref{fig:architecture}, applications submit tasks
through a unified execution interface rather than maintaining separate local,
edge, and cloud execution branches. Context observations and cost estimates are
provided to the MAPE-K controller, whose planning phase selects an execution
tier using either the deterministic policy or the learned Random Forest policy.
The selected plan is subsequently executed through a common dispatch layer.
When utilization at the edge tier exceeds the configured thresholds, the
server-side broker may forward the request to the cloud tier without exposing
this forwarding decision to the mobile application.


\subsection{MAPE-K Adaptation Loop}
\label{subsec:mapek_implementation}

The MAPE-K control cycle is implemented explicitly within the middleware.
The \emph{Monitor} phase acquires the current context state, the
\emph{Analyze} phase derives latency and energy estimates, the \emph{Plan}
phase invokes the active decision policy, the \emph{Execute} phase dispatches
the task to the selected tier, and the \emph{Knowledge} component stores
recent context observations and execution outcomes for subsequent decisions
and empirical analysis.

\paragraph{Monitor.}
Context information is sampled every 2000~ms from collectors covering network
state, processor utilization, battery state, location, and an
accelerometer-derived mobility signal. The resulting observations are combined
into a coherent context snapshot and transformed into normalized decision
features before policy evaluation.

A principal decision variable is the network score. It combines network type,
signal strength, and bandwidth capability using weights of 0.4, 0.3, and 0.3,
respectively, and is subsequently adjusted by an RTT-dependent link-health
factor. The factor equals 1.0 for RTT values below 80~ms and decreases linearly
to zero at 500~ms. The resulting metric therefore incorporates both nominal
link capability and measured communication responsiveness.

\paragraph{Analyze.}
During the analysis phase the middleware computes, for the current context and
task, an estimated latency and energy for both local and remote execution;
these are the quantitative inputs to placement. Task complexity is treated as an
ordinal category, with echo and SHA-256 classified as lightweight, grayscale
conversion as moderate, and matrix multiplication and video-frame edge detection
as computationally intensive, and with per-category baseline compute times of
50, 300, and 2000~ms.

Estimated local latency is the complexity baseline divided by a processor
headroom factor. Estimated remote latency is the sum of four terms: the measured
RTT, an uplink transmission time equal to the input size divided by the
effective bandwidth, a server compute time of 30\% of the local baseline, a
fixed 30~ms queue allowance, and a mobility penalty of up to 200~ms scaled by
the mobility score. When RTT or bandwidth has not yet been measured,
per-network-type defaults are substituted. The energy model mirrors this
structure: estimated local energy is a fixed processor power coefficient
(800~mW) multiplied by the local compute time, while estimated remote energy is
the radio transmit power (1500~mW) multiplied by the transmission time plus a
connected-idle radio power (50~mW) multiplied by the wait time (RTT, server
compute, and queue). The power coefficients are order-of-magnitude figures
intended for per-device calibration. The remote energy term prices only the
radio's transmit and idle draw and does not scale whole-device background
consumption with the full round-trip duration; the consequences of this
simplification, and of the uncalibrated coefficients, are examined in
Section~\ref{sec:cost-validation}.

\paragraph{Plan.}
A planning cycle is initiated synchronously for each submitted task. In
parallel, the middleware performs periodic context-drift detection. Every
3000~ms, context values averaged over a 5000~ms history window are evaluated,
and a drift event is recorded whenever any monitored score changes by at
least 0.2.

The periodic mechanism records environmental change but does not alter the
target of a task already in execution. Updated context therefore affects
subsequent submissions, preserving a clear correspondence between each
task-level placement decision and the context on which it was based.

\paragraph{Execute.}
Tasks are dispatched to the selected tier under a 10~s timeout. If remote
execution fails or exceeds this limit, the middleware re-executes the task
locally and records the fallback event. This recovery mechanism is disabled
for the cloud-only experimental baseline so that remote failures are not
masked by local execution.

\paragraph{Knowledge.}
The knowledge component retains recent contextual observations for drift
detection and stores the inputs, placement decisions, and measured execution
outcomes associated with each task. These records constitute the primary
empirical dataset used in Section~\ref{sec:evaluation}.


\subsection{Deterministic Offloading Policy}
\label{subsec:rule_policy}

The deterministic planning mechanism applies seven rules in a strict
first-match order, prioritizing connectivity and feasibility constraints before
performance-oriented decisions.

Execution remains local when connectivity is unavailable or the network score
falls below 0.30. If the estimated end-to-end speedup is below
$1.5\times$, offloading is also rejected to avoid marginal gains that may not
justify communication overhead. Otherwise, lightweight tasks are directed to
remote execution. A battery-aware rule also permits offloading when the battery
level is below 30\%, the device is not charging, and the estimated remote
energy cost is lower than the local estimate. Computationally intensive tasks
are offloaded when the network score is at least 0.60. Remaining cases are
resolved by an equal-weight latency-energy cost function with a 5\%
hysteresis margin.

When remote execution is selected, the edge tier is preferred if the device is
stationary and the network score is at least 0.60; otherwise, the cloud tier is
used. Local-only and cloud-only modes bypass the adaptive policy and serve as
experimental baselines.

\begin{algorithm}[t]
\caption{Deterministic task-placement procedure used during the planning phase.}
\label{alg:rule-policy}
\begin{algorithmic}[1]
\Require task $\tau$, complexity $c$, context $\phi$, latency estimates $T_l,T_r$, energy estimates $E_l,E_r$
\Ensure execution tier
\State $s \gets T_l/\max(T_r,\epsilon)$
\If{network unavailable \textbf{or} network score $<0.30$}
    \State \textbf{return} local
\EndIf
\If{$s<1.5$}
    \State \textbf{return} local
\EndIf
\If{$c$ is lightweight}
    \State $target \gets \Call{SelectRemoteTier}{\phi}$
\ElsIf{battery $<30\%$ \textbf{and} not charging \textbf{and} $E_r<E_l$}
    \State $target \gets \Call{SelectRemoteTier}{\phi}$
\ElsIf{$c$ is computationally intensive \textbf{and} network score $\geq0.60$}
    \State $target \gets \Call{SelectRemoteTier}{\phi}$
\Else
    \State $J_l \gets 0.5T_l+0.5E_l$
    \State $J_r \gets 0.5T_r+0.5E_r$
    \If{$J_l\leq1.05J_r$}
        \State $target \gets$ local
    \Else
        \State $target \gets \Call{SelectRemoteTier}{\phi}$
    \EndIf
\EndIf
\State \textbf{return} $target$
\Function{SelectRemoteTier}{$\phi$}
    \If{stationary \textbf{and} network score $\geq0.60$}
        \State \textbf{return} edge
    \Else
        \State \textbf{return} cloud
    \EndIf
\EndFunction
\end{algorithmic}
\end{algorithm}

Algorithm~\ref{alg:rule-policy} summarizes the deterministic planning procedure.
The relative order of the latency-sensitive, battery-aware, and compute-oriented
rules is semantically significant because these conditions are evaluated
sequentially rather than independently. Consequently, a lightweight task that
also satisfies the low-battery condition is resolved by the latency-sensitive
rule before the battery-aware rule is considered. This behavior preserves the
first-match semantics of the deployed policy.

The thresholds used in Algorithm~\ref{alg:rule-policy} were fixed at design
time as part of the policy specification rather than derived through a formal
optimization procedure. Their initial values were selected using domain
reasoning and refined during a short calibration phase on the deployed
testbed. Where possible, their behavior is additionally examined against the
collected measurements in Section~\ref{sec:evaluation}. All thresholds are
implemented as configurable policy parameters; systematic sensitivity analysis
and automatic tuning from execution traces are left for future work.

\begin{table}[t]
\centering
\caption{Decision thresholds, rationale, and empirical basis}
\label{tab:thresholds}
\footnotesize
\begin{tabularx}{\columnwidth}{p{2.0cm}cX}
\toprule
Threshold & Value & Rationale and empirical basis \\
\midrule

Unstable network
& $0.30$
& The network score is strongly bimodal in the collected data: 1718 observations
form a healthy cluster (median RTT 34~ms) at score $\geq 0.60$, and 135 form a
degraded cluster (median RTT \textbf{1187~ms}, 90th percentile 1546~ms) at score
$<0.30$, with only 25 of the 1878 measured observations lying in between. The
threshold sits in this sparse gap, making the classification insensitive to its
exact value. Operationally, a median round-trip above 1~s exceeds the local
execution time of every workload, so offloading cannot reduce latency below the
threshold; all 135 sub-threshold observations were executed locally. \\

\addlinespace

Good network / edge eligibility
& $0.60$
& The upper edge of the same bimodal gap: the point above which the healthy
cluster begins (1718 observations, median RTT 34~ms, 90th percentile 74~ms,
within the range where the link-health factor is near unity). Edge selection
and compute-intensive offloading are thereby restricted to links whose measured
responsiveness is empirically high, and the intermediate 0.30 to 0.60 band
(only 25 observations) forms a transition region that avoids abrupt switching. \\

\addlinespace

Compute floor
& $1.5\times$
& Motivated by two measured effects. First, offloading is empirically far
slower than local execution for the compute workloads: median edge latency
exceeded median local latency roughly six-fold (216~ms versus 34~ms for
SHA-256). Second, the remote-latency estimator is only moderately accurate and
its error is tier-dependent (Section~\ref{sec:cost-validation}), so an estimated
speedup near unity cannot be trusted to reflect a real speedup. A margin above
the $1.0\times$ break-even point is therefore required for offloading to be
beneficial in practice, and $1.5\times$ is a conservative choice. This was the
most frequently triggered rule (728 activations), all at estimated speedups
below $1.5\times$ (median $0.26\times$) and predominantly on lightweight
tasks. \\

\addlinespace

Low battery
& $<30\%$
& A conventional mobile power-management threshold rather than a data-derived
optimum. The discharge sweeps induced levels of 8 to 28\%, activating the rule
154 times; because these levels were forced rather than naturally observed, the
dataset exercises the rule but cannot establish 30\% as optimal. The condition
is further guarded by requiring the estimated remote energy to be below the
local estimate, so it never offloads when doing so would cost more energy. \\

\addlinespace

Hysteresis margin
& $5\%$
& A standard dead-band that suppresses decision oscillation when the estimated
local and remote costs are close. In the collected data all 76 balanced-cost
decisions fell far from the boundary (none within $\pm 15\%$ of it), so the
margin did not alter any recorded decision; its specific value is thus a design
safeguard against oscillation rather than a quantity the present data can
calibrate. \\

\bottomrule
\end{tabularx}
\end{table}


\subsection{Random Forest Policy}
\label{subsec:rf_policy}

The learned planning backend is implemented using a Random Forest classifier.
The execution mode selects which planning backend is active for a given task:
the deterministic policy, the learned policy, or one of the two static
experimental baselines.

The deployed classifier contains 50 decision trees. After training with
Scikit-learn, each tree is serialized to JSON using its feature indices,
threshold values, child-node indices, and class probabilities at terminal
nodes. The Android middleware evaluates this representation using a lightweight
custom tree interpreter, avoiding dependence on a dedicated mobile
machine-learning runtime.

Each tree is traversed from its root until a terminal node is reached. The
class-probability vectors produced by all trees are accumulated and normalized.
This procedure reproduces the soft-voting behavior of the original
Scikit-learn classifier, and the resulting class is mapped to local, edge, or
cloud execution.


\begin{algorithm}[t]
\caption{On-device inference procedure used by the Random Forest planning policy.}
\label{alg:rf-policy}
\begin{algorithmic}[1]
\Require task $\tau$, context $\phi$, forest of $N=50$ trees
\Ensure execution tier
\State $x \gets \Call{ExtractFeatures}{\tau,\phi}$
\State $v \gets \mathbf{0}$
\For{each tree $t$}
    \State $n \gets 0$
    \State $k \gets 0$
    \While{$n$ is not a leaf}
        \If{$x[f_n]\leq\theta_n$}
            \State $n \gets left_n$
        \Else
            \State $n \gets right_n$
        \EndIf
        \State $k \gets k+1$
        \If{$k$ exceeds tree size}
            \State \textbf{abort} malformed model
        \EndIf
    \EndWhile
    \State $v \gets v+p_n$
\EndFor
\State $p \gets v/\sum v$
\State $y \gets \operatorname*{arg\,max}(p)$
\If{$y$ represents edge}
    \State $target \gets$ edge
\ElsIf{$y$ represents cloud}
    \State $target \gets$ cloud
\Else
    \State $target \gets$ local
\EndIf
\State \textbf{return} $target$
\Function{ExtractFeatures}{$\tau,\phi$}
    \State \textbf{return} $\langle b,q,n,s_n,r,w,u,m,c,T_l,T_r,T_l/\max(T_r,\epsilon)\rangle$
\EndFunction
\end{algorithmic}
\end{algorithm}

Algorithm~\ref{alg:rf-policy} uses twelve contextual and estimated-cost features: battery level,
charging state, network type, network score, RTT, bandwidth capability,
processor utilization, mobility stability, task complexity, estimated local
latency, estimated remote latency, and expected speedup.

The feature order is fixed and validated against the metadata exported during
training. This check prevents silent inconsistencies between the training and
deployment pipelines caused by reordered or renamed variables. On-device
inference required less than 1~ms in the deployed implementation.

To preserve decision availability, failure to load or evaluate the learned
model triggers the deterministic policy. Consequently, the learned execution
mode does not make task placement dependent on successful model inference and
retains a deterministic fallback mechanism.

Because the learned policy is trained on the deterministic policy's own
decisions, it is an on-device imitation of the rule engine rather than an
independent optimizer. Its purpose here is to establish that a compact learned
model can reproduce the rule engine's placements at negligible inference cost;
its runtime decision quality is quantified against the rule engine in
Section~\ref{sec:decision-quality}, where it is found to be at least on par
rather than superior.


\subsection{Communication and Server-Side Execution}
\label{subsec:communication_server}

Reachability and RTT for the remote tiers are obtained through lightweight
health probes. Results are cached for 5~s to avoid introducing an additional
network request for every task submission. Before the first probe completes,
a zero-valued stale state is used so that freshness calculations remain
numerically well defined.

Before remote dispatch, the middleware consults the cached reachability state.
An endpoint already known to be unavailable is rejected immediately rather
than incurring the complete connection timeout. For reachable endpoints, the
task payload is encoded using Base64 and submitted to the offloading service.
The Android client uses MIME-compatible Base64 decoding to accommodate
line-wrapped server responses.

On the server side, binary payload fields are represented using a Base64-aware
Pydantic data type so that encoded JSON strings are decoded into their original
binary representation rather than interpreted as ordinary UTF-8 content.

Both the edge and cloud services dispatch requests through a common task
registry containing handlers for echo, SHA-256, grayscale conversion, matrix
multiplication, and video-frame edge detection. Consequently, equivalent
remote workloads use the same registered implementation on both tiers rather
than independently maintained copies.

The edge service additionally limits executor concurrency to four workers and
reports queue-state information. Processor and memory utilization are measured
relative to the container resource limits. When processor utilization exceeds
85\% or memory utilization exceeds 80\%, the edge broker forwards the original
request to the cloud service and relays the resulting response back to the
client. This server-side adaptation is transparent to the mobile device, which
continues to interact with the originally selected edge endpoint and receives no
explicit redirection.


\subsection{Task Abstraction and Instrumentation}
\label{subsec:task_instrumentation}

Each offloadable unit of computation is represented by a task abstraction
containing an identifier, task name, input size, complexity category, input
payload, and local execution function. Workloads are categorized as
lightweight, moderate, or computationally intensive according to their
processing characteristics.

For remote execution, the middleware serializes the task name, complexity
category, and input payload into an HTTP request, after which the corresponding
server-side handler performs the operation. This local and remote execution
symmetry is maintained for four of the five benchmark workloads.

The video-frame edge-detection workload is an explicit exception. Because
OpenCV is unavailable in the deployed Android implementation, the local
fallback extracts the first video frame as a JPEG rather than applying the
Canny edge-detection operation performed on the remote tiers. This limitation
is reported explicitly because the local and remote execution paths are not
computationally equivalent for this workload.

The instrumentation layer records one observation for every task execution.
Each record contains the policy outcome, selected and actual execution tiers,
fallback status, measured latency and energy, and the contextual variables
available to the decision policy. Measured energy is obtained from the platform
battery-current property (BATTERY\_PROPERTY\_CURRENT\_NOW) sampled immediately
before and after each execution, integrated as whole-device power over the
execution window; it is therefore an upper bound that also includes screen and
background draw, and is available only on devices that implement the property.
This structure allows the same execution trace to support latency, reliability,
estimator, energy, and decision-quality analyses.

Execution records additionally identify experiments in which a contextual
variable was deliberately controlled to exercise infrequent policy branches.
These observations can therefore be retained for rule-coverage and classifier
training while being excluded from evaluations that require naturally measured
context, such as estimator validation.

Overall, the implementation supports local, edge, and cloud execution through
actual task dispatch and HTTP communication while preserving a common
monitoring, communication, execution, fallback, and instrumentation pipeline.
Accordingly, differences between the deterministic and learned adaptive modes
can be attributed primarily to the planning mechanism rather than to
differences in the surrounding execution infrastructure.
